\section{Conclusion}
We integrated planning, alternative retrieval paths, cited generation, and bounded revision. The native fine-tuned preset directly ranks Dense evidence without reranking, while hybrid search remains a comparator. Its retrieval study favors struct passages and shows a promising adaptation--latency tradeoff within a development protocol. Historical hybrid downstream results support domain-aware context allocation, while independent generation adds cost without a final coverage advantage. The loop yields modest acceptance gains and a tradeoff between withholding drafts and preserving partial information. Conditional errors and unstable grades motivate matched, independently annotated evaluation.

\section*{Limitations}
Our results are development evidence from one dataset and largely one model configuration. The synthetic requests have explicit domain cues, provisional annotations, and unverified alternative evidence sources. The retrieval comparison uses a distributed dummy-file cohort; its gains and CUDA timings should not be generalized to downstream generation or other hardware. The official validation set has already informed development. Human factuality, hallucination, user satisfaction, completed task execution, and successful human handoff are unmeasured. Stored experiments predate the namespace refactor and subsequent wrapper hardening; those changes have not been re-evaluated on these cohorts.
